\documentclass[11pt,twoside]{amsart}
\usepackage{amssymb, amsmath, enumerate, palatino, hyperref}
\usepackage[normalem]{ulem}
\usepackage{fullpage}
\usepackage[T1]{fontenc}
\renewcommand{\labelitemi}{\guillemotright}
\usepackage{mathrsfs}


\theoremstyle{plain}
\newtheorem{prop}{Proposition}%[section]
\newtheorem{lemma}[prop]{Lemma}
\newtheorem{thm}[prop]{Theorem}
\newtheorem{obs}[prop]{Observation}
\newtheorem{app}[prop]{Application}
\newtheorem*{MainThm}{Main Theorem}
\newtheorem{cor}[prop]{Corollary}
\newtheorem{conj}[prop]{Conjecture}
\theoremstyle{remark}
\newtheorem{rmk}[prop]{Remark}
\newtheorem{prob}[prop]{Problem}
\newtheorem{q}[prop]{Question}
\newtheorem{bonus}[prop]{Bonus Problem}
\newtheorem{exc}{Exercise}
\theoremstyle{definition}
\newtheorem{ex}[prop]{Example}
\theoremstyle{definition}
\newtheorem{defn}[prop]{Definition}

\newcommand{\RR}{\mathbb{R}}
\newcommand{\ZZ}{\mathbb{Z}}
\newcommand{\CC}{\mathbb{C}}
\newcommand{\NN}{\mathbb{N}}
\newcommand{\QQ}{\mathbb{Q}}
\newcommand{\PP}{\mathbb{P}}

\newcommand{\cC}{\mathscr{C}}
\newcommand{\Ob}{\operatorname{Ob}}
\newcommand{\Mor}{\operatorname{Mor}}

\newcommand{\Set}{\operatorname{Set}}
\newcommand{\FinSet}{\operatorname{FinSet}}
\newcommand{\Vect}{\operatorname{Vect}}
\newcommand{\FinVect}{\operatorname{FinVect}}
\newcommand{\Mat}{\operatorname{Mat}}
\newcommand{\Gp}{\operatorname{Gp}}
\newcommand{\FinGp}{\operatorname{FinGp}}
\newcommand{\AbGp}{\operatorname{AbGp}}
\newcommand{\Ring}{\operatorname{Ring}}
\newcommand{\CommRing}{\operatorname{CommRing}}
\newcommand{\Field}{\operatorname{Field}}
\newcommand{\Top}{\operatorname{Top}}

\newcommand{\Aut}{\operatorname{Aut}}

\newcommand{\defeq}{:=}

\title{Math 113: Discrete Structures\\ Homework 27}

\begin{document}
\maketitle




\begin{prob}
What is the expected value of the number of digits equal to 3 in a 4-digit
positive integer? Write your solution as a fraction~$a/b$ in lowest terms.  The sample space is
\[
  S = \left\{ a_1a_2a_3a_4: a_1\in\left\{ 1,2,\dots,9 \right\},
  a_2,a_3,a_4\in\left\{ 0,1,\dots9 \right\} \right\}
\]
 [Hint: express the relevant random variable
as a sum of simpler random variables, and use linearity of expectation.]
\end{prob}

\begin{prob}
Let $\pi$ be a permutation of $\underline{n}$.  The index $i$ is called an
\emph{exceedance} of $\pi$ if $\pi(i)>i$.  For instance, using the
notation~$\pi(1),\pi(2),\dots,\pi(n)$ for a permutation~$\pi$, the
permutation~$\pi=3,2,4,1$ has exceedance~$2$ since~$\pi(1)=3>1$
and~$\pi(3)=4>3$ (but $\pi(2)=2$ and $\pi(4)=1<4$).  
\begin{enumerate}[(a)]
  \item Let~$X_i$ be the random variable on the set of permutations
    of~$\underline{n}$ such that~$X_i(\pi)=1$ if~$i$ is an exceedance
    of~$\pi$, and~$X_i(\pi)=0$, otherwise.  What is the expected value,~$E(X_i)$?
  \item How many exceedances does the average permutation
    of~$\underline{n}$ have?
\end{enumerate}
\end{prob}

\end{document}
